\begin{sectionbox}
    \subsection{Allgemeines}
    \textbf{Allgemeine Wellengleichung:} $\frac{1}{c^2}\frac{\partial^2u}{\partial t^2}=\Delta u$\\

    $c=\lambda   f=\frac{\omega}{k}=\frac{x}{\Delta t}$\\
    $k=\frac{2\pi}{\lambda}=\frac{\omega}{c} \qquad \text{Wellenzahl}$ \\

    \subsubsection{Laufende Welle}
    \begin{emphbox}
        \textbf{Funktion:} $y(x,t)=y_{\ir max}\sin(kx \cpm \omega t + \varphi_0)^{\textcolor{red}{\text{+: Welle in neg. x}}}_{\textcolor{green}{\text{-: Welle in pos. x}}}$
    \end{emphbox}

    \subsubsection{Stehende Welle}
    Oberschwingung: $n$, ($n=0$ für $0$-te Schwingung $\rightarrow$ ein Bauch) \\
    Stablänge: $L=\frac{(2n+1)\lambda}{2}$ \\
    Grundfrequenz: $f_0=\frac{c}{2L}$ \\
    Frequenzen der Oberschwingungen: $f_n=(2n+1)f_0$
\end{sectionbox}

\begin{sectionbox}
    \subsection{Doppler-Effekt}
    \begin{cookbox}{}
        \item Beobachter \textbf{bewegt}, Quelle \textbf{ruht}\\
        \begin{itemize}
            \item[] $f_B=f_Q(1\cpm\frac{v_B}{c})^{\textcolor{red}{\text{+: Annäherung}}}_{\textcolor{green}{\text{-: Entfernung}}}$\\
        \end{itemize}
        \item Beobachter \textbf{ruht}, Quelle \textbf{bewegt}\\
        \begin{itemize}
            \item[] $f_B=\frac{f_Q}{1\cpm v_Q/c}^{\textcolor{red}{\text{+: Entfernung}}}_{\textcolor{green}{\text{-: Annäherung}}}$\\
        \end{itemize}
        \item Beobachter \textbf{bewegt}, Quelle \textbf{bewegt}\\
        $f_B = f_Q \frac{c \, \underset{\textcolor{green}{-}}{\textcolor{red}{+}} \, V_B}{c \, \underset{\textcolor{green}{+}}{\textcolor{red}{-}} \, V_Q} \qquad \text{Q} \, \underset{\textcolor{green}{\leftarrow} \textcolor{green}{\rightarrow}}{\textcolor{red}{\rightarrow} \textcolor{red}{\leftarrow}} \, \text{B}$\\
        $f_B = f_Q \frac{c \, \underset{\textcolor{cyan}{-}}{\textcolor{purple}{+}} \, V_B}{c \, \underset{\textcolor{cyan}{-}}{\textcolor{purple}{+}} \, V_Q} \qquad \text{Q} \, \underset{\textcolor{cyan}{\rightarrow} \textcolor{cyan}{\rightarrow}}{\textcolor{purple}{\leftarrow} \textcolor{purple}{\leftarrow}} \, \text{B}$\\
        \item Quelle \textbf{bewegt} (TODO) ($v>1 \, \ir Ma$)\\
        $\sin \alpha=\frac{\diff w}{\diff s}=\frac{c_s t}{v_Qt}=\frac{1}{\ir Ma}$
    \end{cookbox}
\end{sectionbox}


\begin{sectionbox}
    \subsection{Interferenz}
    \subsubsection{Einfachspalt (Beugung)}
    $b$: Spaltbreite \\
    Gangunterschied: $\Delta s = b\sin\alpha \approx \frac{x}{D}b$ \\
    Maxima: $\Delta s=\pm (m + \frac{1}{2})\lambda$ \\
    Minima: $\Delta s=\pm m\lambda \rightarrow x_{\ir min} \approx \pm\frac{m\lambda D}{b} \qquad m = (1,2,...,n)$ \\

    \subsubsection{Doppelspalt}
    $d$: Spaltenabstand, $D$: Abstand zum Schirm, $m = (0,1,...,n)$ \\
    Gangunterschied: $\Delta s = d\sin\alpha \approx \frac{x}{D}d$ \\
    Phasenunterschied: $\Delta\varphi = 2\pi \frac{\Delta s}{\lambda}$ \\
    Konstr. Int. (Maxima): $\Delta s=\pm m\lambda \rightarrow x_{\ir max}\approx\pm\frac{m\lambda D}{d}$ \\
    Destr. Int. (Minima): $\Delta s = \pm (m - \frac{1}{2})\lambda \rightarrow x_{\ir min}\approx \pm (m - \frac{1}{2})\frac{\lambda D}{d}$ \\

    \subsubsection{Gitter}
    $d$: Spaltenabstand, $D$: Abstand zum Schirm, $m = (0,1,...,n)$ \\
    Gangunterschied: $\Delta s = d\sin\alpha \approx \frac{x}{D}d$ \\
    Hauptmaxima: $\Delta s=\pm m\lambda \rightarrow x_{\ir max}\approx\pm\frac{m\lambda D}{d}$ \\
    Nebenmaxima: $\Delta s = \pm \frac{2k + 1}{2N}\lambda$\\
    Abstand benachbarter Maxima $\Delta x \approx \frac{\lambda D}{d}$ \\

    \subsubsection{An dünner Schicht $n_2 > n_1$}
    $d$: Dicke der Schicht, $\alpha$: Einfallswinkel \\
    $\Delta s = 2d \sqrt{n_2^2-n_1^2\sin^2\alpha} + \frac{\lambda}{2}$ \\
    Maxima und Minima wie beim Doppelspalt
\end{sectionbox}